
\section{\BookRecordIndex{T0141}格罗莫夫收敛}\label{sec:5-2}
\zhindex{geluomofushoulian}{格罗莫夫收敛}

在本章其余部分，我们将发展\BookRecordIndex{T0340}稳定映射列\BookRecordIndex{T0141}格罗莫夫收敛的概念。它是第 D.5 节中\BookRecordIndex{T0339}稳定曲线
相应概念的细化；\BookRecordIndex{T0339}稳定曲线的情形不涉及分析，因而简单得多。

下面正式定义一列 $J^\nu$-全纯曲线 $u^\nu:S^2\to M$ 收敛到\BookRecordIndex{T0340}稳定映射
$(\mathbf u,\mathbf z)\in\mathcal{SC}_{0,n}(M;J)$ 的含义。回顾第 D.1 节：若映射列
$v^\nu:S^2\to M$ 在 $S^2\setminus X$ 的每个紧子集上按 $C^\infty$ 拓扑收敛到
$v:S^2\to M$，就称它在 $S^2\setminus X$ 上\emph{\BookRecordIndex{T0416}逐紧一致收敛}（u.c.s.）。还要回顾
定义~\ref{def:5-1-1} 中的记号
\[
  Z_\alpha:=\{z_{\alpha\beta}\mid \alpha E\beta\},
  \qquad
  Y_\alpha:=Z_\alpha\cup\{z_i\mid \alpha_i=\alpha\}.
\]

\begin{definition}[\BookRecordIndex{T0141}格罗莫夫收敛]\label{def:5-2-1}
设 $J^\nu$ 是 $M$ 上一列被 $\omega$ \BookRecordIndex{T0372}驯服的\BookRecordIndex{T0088}殆复结构，它在 $C^\infty$ 拓扑下收敛到
$J\in\mathcal J_\tau(M,\omega)$；设
\[
  (\mathbf u,\mathbf z)
  =\bigl(\{u_\alpha\}_{\alpha\in T},
    \{z_{\alpha\beta}\}_{\alpha E\beta},
    \{\alpha_i,z_i\}_{1\le i\le n}\bigr)
  \in\mathcal{SC}_{0,n}(M;J)
\]
是一个\BookRecordIndex{T0340}稳定映射；又设 $u^\nu:S^2\to M$ 是一列带有 $n$ 个两两不同\BookRecordIndex{T0066}标记点
$z_1^\nu,\ldots,z_n^\nu\in S^2$ 的 $J^\nu$-全纯映射。若存在一族 \BookRecordIndex{T0040}Möbius 变换
\[
  \{\phi_\alpha^\nu\}_{\alpha\in T,\,\nu\in\mathbb N},
\]
使下列条件成立，就称映射列
$(u^\nu,\mathbf z^\nu):=(u^\nu,z_1^\nu,\ldots,z_n^\nu)$
\emph{\BookRecordIndex{T0141}格罗莫夫收敛}到 $(\mathbf u,\mathbf z)$。

\medskip
\noindent\textup{（映射）}
对每个 $\alpha\in T$，映射列
\[
  u_\alpha^\nu:=u^\nu\circ\phi_\alpha^\nu:S^2\to M
\]
在 $S^2\setminus Z_\alpha$ 上\BookRecordIndex{T0416}逐紧一致收敛到 $u_\alpha$。

\medskip
\noindent\textup{（能量）}
若 $\alpha E\beta$，且 $m_{\alpha\beta}(\mathbf u)$ 由式~\eqref{eq:5-1-1} 定义，则
\begin{equation}\label{eq:5-2-1}
  m_{\alpha\beta}(\mathbf u)
  =\lim_{\varepsilon\to0}\lim_{\nu\to\infty}
  E\bigl(u_\alpha^\nu;B_\varepsilon(z_{\alpha\beta})\bigr).
\end{equation}

\medskip
\noindent\textup{（重标度）}
若 $\alpha E\beta$，则映射列
\[
  \phi_{\alpha\beta}^\nu
  :=(\phi_\alpha^\nu)^{-1}\circ\phi_\beta^\nu
\]
在 $S^2\setminus\{z_{\beta\alpha}\}$ 上\BookRecordIndex{T0416}逐紧一致收敛到 $z_{\alpha\beta}$。

\medskip
\noindent\textup{（\BookRecordIndex{T0066}标记点）}
对 $i=1,\ldots,n$，有
\[
  z_i=\lim_{\nu\to\infty}
  (\phi_{\alpha_i}^\nu)^{-1}(z_i^\nu).
\]
\end{definition}

这个定义既复杂又重要，值得花些篇幅解释。理解它的一处困难在于，原映射列 $u^\nu$
除了出现在复合
\[
  u_\alpha^\nu:=u^\nu\circ\phi_\alpha^\nu
\]
中以外，并不显式出现；这是因为我们只关心它在群 $G$ 的作用模掉之后的极限行为。
因此，最好从某个顶点 $\alpha$ 的视角观察收敛。（映射）公理说明，$u_\alpha^\nu$
在 $S^2\setminus Z_\alpha$ 的所有紧子集
$S^2\setminus B_\varepsilon(Z_\alpha)$ 上收敛。进一步，对每个与 $\alpha$ 相邻的顶点
$\beta$，相应映射 $u_\beta^\nu$ 是 $u_\alpha^\nu$ 的重标度，含义如下。

所谓一个\emph{以 $x$ 为源、以 $y$ 为靶的 $S^2$ 重标度}，是指一列\BookRecordIndex{T0126}分式线性变换
$\phi^\nu\in G$，它在 $S^2\setminus\{x\}$ 上\BookRecordIndex{T0416}逐紧一致收敛到 $y$。设
$\phi_i^\nu$ 对 $i=0,1$ 分别是以 $x_i$ 为源、以 $y_i$ 为靶的重标度列。若
$y_0\ne x_1$，则复合 $\phi_1^\nu\circ\phi_0^\nu$ 是以 $x_0$ 为源、以 $y_1$
为靶的重标度列。此外，由习题 D.1.3，以 $x$ 为源、以 $y$ 为靶的 $S^2$ 重标度列之逆，
是以 $y$ 为源、以 $x$ 为靶的重标度列。用这一术语，（重标度）公理说明
$\phi_{\alpha\beta}^\nu$ 是以 $z_{\beta\alpha}$ 为源、以 $z_{\alpha\beta}$ 为靶的
重标度列。由于当 $\alpha\ne\gamma$ 时
$z_{\beta\alpha}\ne z_{\beta\gamma}$，复合
\[
  \phi_{\alpha\gamma}^\nu
  =\phi_{\alpha\beta}^\nu\circ\phi_{\beta\gamma}^\nu
\]
也是一个重标度列，其源为 $z_{\gamma\beta}$、靶为 $z_{\alpha\beta}$。\BookRecordIndex{T0411}重标度后的映射列
\[
  u_\beta^\nu=u_\alpha^\nu\circ\phi_{\alpha\beta}^\nu
\]
在 $S^2\setminus B_\varepsilon(Z_\beta)$ 的紧子集上收敛。它从
$u_\alpha^\nu$ 在\BookRecordIndex{T0243}冒泡点 $z_{\alpha\beta}$ 的小邻域
\[
  \phi_{\alpha\beta}^\nu
  \bigl(S^2\setminus B_\varepsilon(z_{\beta\alpha})\bigr)
\]
中截取出一部分并将其“放大”。类似地，从 $\beta$-\BookRecordIndex{T0275}气泡的视角观察 $\alpha$-\BookRecordIndex{T0275}气泡时，
$u_\alpha^\nu$ 是 $u_\beta^\nu$ 在 $z_{\beta\alpha}$ 的小邻域
\[
  \phi_{\beta\alpha}^\nu
  \bigl(S^2\setminus B_\varepsilon(z_{\alpha\beta})\bigr)
\]
中那一部分的放大。

（\BookRecordIndex{T0066}标记点）公理说明\BookRecordIndex{T0066}标记点收敛，而（能量）公理说明没有能量损失。为解释后一条件的
重要性，定义映射列 $v^\nu:S^2\to M$ 在点 $z$ 处的\emph{\BookRecordIndex{T0276}气泡能量}为极限
\[
  E(\{v^\nu\};z)
  :=\lim_{\varepsilon\to0}\lim_{\nu\to\infty}
  E\bigl(v^\nu;B_\varepsilon(z)\bigr).
\]
每当使用\BookRecordIndex{T0276}气泡能量时，我们都假设右端极限存在，并且对某个 $\varepsilon>0$，映射列
$v^\nu$ 在 $B_\varepsilon(z)\setminus\{z\}$ 上\BookRecordIndex{T0416}逐紧一致收敛。在这种情形下，称
$v^\nu$ \emph{在 $z$ 处伴随\BookRecordIndex{T0243}冒泡收敛}。依此术语，（能量）公理断言：映射列
$u_\alpha^\nu$ 在 $z_{\alpha\beta}$ 处的\BookRecordIndex{T0276}气泡能量，等于极限\BookRecordIndex{T0340}稳定映射
$(\mathbf u,\mathbf z)$ 在从 $\alpha$ 指向 $\beta$ 的边所分离出的那部分树上的能量
$m_{\alpha\beta}(\mathbf u)$。这意味着重标度的选取没有损失能量。特别地，若
$m_{\alpha\beta}(\mathbf u)>0$，则映射列
\[
  \psi^\nu:=\phi_{\alpha\beta}^\nu
\]
满足命题~\ref{prop:4-7-2} 的要求，其中
\[
  u^\nu=u_\alpha^\nu,
  \qquad v^\nu=u_\beta^\nu,
  \qquad z_0:=z_{\alpha\beta},
  \qquad z_\infty:=z_{\beta\alpha}.
\]

定义~\ref{def:5-2-1} 的（能量）公理蕴含命题~\ref{prop:4-7-1} 的假设 (b)。因此，
它包含了相邻顶点 $\alpha,\beta$ 所对应的\BookRecordIndex{T0275}气泡彼此相连这一信息。于是，只要建立这种形式
的能量估计，就能证明格罗莫夫极限是连通的。

后面的各节将证明：极限在等价意义下唯一，而且每一列能量一致有界的
$J^\nu$-全纯球面 $u^\nu$ 都有一个在定义~\ref{def:5-2-1} 意义下收敛的子列。
作为开端，我们先\BookRecordIndex{T0326}推出\BookRecordIndex{T0139}格罗莫夫紧性的几个并非完全直接的后果。例如，若
$x^\nu\in u^\nu(S^2)$ 收敛，则其极限落在极限曲线的像中；同伦类在\BookRecordIndex{T0141}格罗莫夫收敛下
也保持不变。注意，下述 (iii) 中的点 $z_{\alpha i}$ 由式~\eqref{eq:5-1-6} 定义。
此外，(iv) 关于\BookRecordIndex{T0185}交比 $w_{ijk\ell}$ 的断言说明：\BookRecordIndex{T0141}格罗莫夫收敛列在\BookRecordIndex{T0377}遗忘映射下的像，
在\BookRecordIndex{T0339}稳定曲线\BookRecordIndex{T0246}模空间中收敛（见附录 D）。

\begin{theorem}\label{thm:5-2-2}
设 $(u^\nu,\mathbf z^\nu)$ 如定义~\ref{def:5-2-1} 所述，\BookRecordIndex{T0141}格罗莫夫收敛到\BookRecordIndex{T0340}稳定映射
$(\mathbf u,\mathbf z)$。则：
\begin{enumerate}[label=(\roman*),leftmargin=3em]
\item 若 $x^\nu\in u^\nu(S^2)$ 收敛到 $x\in M$，则对某个 $\alpha\in T$，有
$x\in u_\alpha(S^2)$；
\item 当 $\nu$ 充分大时，映射 $u^\nu:S^2\to M$ 与\BookRecordIndex{T0227}连通和
$\#_{\alpha\in T}u_\alpha$ 同伦；
\item 对每个 $\alpha\in T$ 和每个 $i$，
\[
  z_{\alpha i}=\lim_{\nu\to\infty}
  (\phi_\alpha^\nu)^{-1}(z_i^\nu);
\]
\item 对所有 $i,j,k,\ell$，
\[
  w_{ijk\ell}(\mathbf z)
  =\lim_{\nu\to\infty}w_{ijk\ell}(\mathbf z^\nu).
\]
\end{enumerate}
\end{theorem}

\begin{proof}
先证 (i)。设 $x^\nu=u^\nu(z^\nu)$ 收敛到 $x$。必要时取子列，不失一般性，可假设
对每个 $\alpha$，点列 $(\phi_\alpha^\nu)^{-1}(z^\nu)$ 都收敛到某个
$z_\alpha\in S^2$。若存在 $\alpha\in T$，使得对每个满足 $\alpha E\beta$ 的
$\beta\in T$ 都有 $z_\alpha\ne z_{\alpha\beta}$，则
\[
  x=\lim_{\nu\to\infty}u^\nu(z^\nu)
   =\lim_{\nu\to\infty}
    u_\alpha^\nu\bigl((\phi_\alpha^\nu)^{-1}(z^\nu)\bigr)
   =u_\alpha(z_\alpha),
\]
结论成立。

若不存在这样的 $\alpha$，我们断言树 $T$ 中存在一条边 $\alpha E\beta$，使得
\begin{equation}\label{eq:5-2-2}
  z_\alpha=z_{\alpha\beta},
  \qquad z_\beta=z_{\beta\alpha}.
\end{equation}
事实上，从任意顶点 $\alpha_0$ 出发，存在唯一的顶点列
$\alpha_0,\alpha_1,\alpha_2,\ldots$，使得
\[
  \alpha_iE\alpha_{i+1},
  \qquad z_{\alpha_i}=z_{\alpha_i\alpha_{i+1}}.
\]
必有某个 $j$ 满足 $\alpha_j=\alpha_{j+2}$，否则这些 $\alpha_i$ 将构成两两不同顶点的
无限序列。于是 $\alpha=\alpha_j$ 与 $\beta=\alpha_{j+1}$ 满足
式~\eqref{eq:5-2-2}。

若 $m_{\alpha\beta}(\mathbf u)=0$，则由式~\eqref{eq:5-2-1} 和
定理~\ref{thm:4-6-1}，映射列 $u_\alpha^\nu$ 在 $z_{\alpha\beta}$ 的一个邻域上一致
收敛到 $u_\alpha$。故仍有
\[
  x=\lim_{\nu\to\infty}u^\nu(z^\nu)=u_\alpha(z_{\alpha\beta}).
\]
若 $m_{\alpha\beta}(\mathbf u)\ne0$，则命题~\ref{prop:4-7-2} 可用于映射列
$u_\alpha^\nu$，其中
\[
  u:=u_\alpha,
  \qquad v:=u_\beta,
  \qquad z_0:=z_{\alpha\beta},
  \qquad z_\infty:=z_{\beta\alpha},
  \qquad \psi^\nu:=\phi_{\alpha\beta}^\nu.
\]
由式~\eqref{eq:5-2-2}，$z_0=z_\alpha$ 且 $z_\infty=z_\beta$，所以
\[
  z_0=\lim_{\nu\to\infty}\zeta^\nu,
  \qquad
  z_\infty=\lim_{\nu\to\infty}(\psi^\nu)^{-1}(\zeta^\nu),
  \qquad
  \zeta^\nu:=(\phi_\alpha^\nu)^{-1}(z^\nu).
\]
命题~\ref{prop:4-7-2} 因而说明
$u^\nu(z^\nu)=u_\alpha^\nu(\zeta^\nu)$ 收敛到
$u(z_0)=v(z_\infty)$
\sourcefootnote{原书此处写作 $u(z_0)=v(\infty)$；但本段已定义
$z_\infty:=z_{\beta\alpha}$，命题~\ref{prop:4-7-2} 的坐标不变结论及紧接着的
$u_\beta(z_{\beta\alpha})$ 都唯一要求此处为 $v(z_\infty)$。}。于是
\[
  x=\lim_{\nu\to\infty}u^\nu(z^\nu)
   =u_\alpha(z_{\alpha\beta})
   =u_\beta(z_{\beta\alpha}),
\]
这就证明了 (i)。

下面证明 (ii)。固定一条边 $\alpha E\beta$。对小于 $2$-球直径的 $r>0$ 和
$\nu\in\mathbb N$，定义
\[
  \begin{aligned}
  A_r^\nu(\alpha,\beta)
  &:=\phi_\alpha^\nu\bigl(B_r(z_{\alpha\beta})\bigr)
    \cap\phi_\beta^\nu\bigl(B_r(z_{\beta\alpha})\bigr)\\
  &=\bigl\{z\in S^2\bigm|
    d\bigl((\phi_\alpha^\nu)^{-1}(z),z_{\alpha\beta}\bigr)<r,\ 
    d\bigl((\phi_\beta^\nu)^{-1}(z),z_{\beta\alpha}\bigr)<r\bigr\}.
  \end{aligned}
\]
因为 $\phi_{\beta\alpha}^\nu$ 在 $S^2\setminus\{z_{\alpha\beta}\}$ 上\BookRecordIndex{T0416}逐紧一致收敛到
$z_{\beta\alpha}$，只要 $\nu\ge\nu_0(r)$ 充分大，像
\[
  \phi_{\beta\alpha}^\nu
  \bigl(S^2\setminus B_r(z_{\alpha\beta})\bigr)
\]
就是 $B_r(z_{\beta\alpha})$ 中的一个小圆盘。因此，当 $\nu\ge\nu_0(r)$ 时，集合
$A_r^\nu(\alpha,\beta)$ 与
\[
  \phi_{\beta\alpha}^\nu\bigl(B_r(z_{\alpha\beta})\bigr)
  \cap B_r(z_{\beta\alpha})
\]
微分同胚，从而是一个\BookRecordIndex{T0170}环域。我们将证明：对每个 $\varepsilon>0$，都存在常数
$r>0$ 和 $\nu_0\in\mathbb N$，使得
\begin{equation}\label{eq:5-2-3}
  \nu\ge\nu_0
  \quad\Longrightarrow\quad
  \sup_{z\in A_r^\nu(\alpha,\beta)}
  d\bigl(u^\nu(z),u_\alpha(z_{\alpha\beta})\bigr)<\varepsilon.
\end{equation}

先设 $m_{\alpha\beta}(\mathbf u)=0$。此时在 $z_{\alpha\beta}$ 处没有\BookRecordIndex{T0243}冒泡，所以
$u_\alpha^\nu$ 在 $z_{\alpha\beta}$ 的一个邻域上一致收敛到 $u_\alpha$。给定
$\varepsilon>0$，取充分小的 $r>0$，使得对
$\zeta\in B_r(z_{\alpha\beta})$ 有
\[
  d\bigl(u_\alpha(\zeta),u_\alpha(z_{\alpha\beta})\bigr)<\varepsilon/2,
\]
并且 $u_\alpha^\nu$ 在 $B_r(z_{\alpha\beta})$ 上一致收敛到 $u_\alpha$。再取
$\nu_0$，使得当 $\zeta\in B_r(z_{\alpha\beta})$ 且 $\nu\ge\nu_0$ 时，
\[
  d\bigl(u_\alpha^\nu(\zeta),u_\alpha(\zeta)\bigr)<\varepsilon/2.
\]
由三角不等式，式~\eqref{eq:5-2-3} 成立。

再设 $m_{\alpha\beta}(\mathbf u)>0$。命题~\ref{prop:4-7-2} 可用于映射列
$u^\nu:=u_\alpha^\nu$、
$\psi^\nu:=\phi_{\alpha\beta}^\nu$
\sourcefootnote{原书此处写作 $\psi^\nu:=\phi_{\beta\alpha}^\nu$；但按
$\phi_{\alpha\beta}^\nu=(\phi_\alpha^\nu)^{-1}\circ\phi_\beta^\nu$，只有取
$\psi^\nu=\phi_{\alpha\beta}^\nu$ 才同时满足
$u_\beta^\nu=u_\alpha^\nu\circ\psi^\nu$，并使下文
$(\psi^\nu)^{-1}((\phi_\alpha^\nu)^{-1}(z))
=(\phi_\beta^\nu)^{-1}(z)$。原书的反向次序不能\BookRecordIndex{T0326}推出所需距离估计。}，映射
$u:=u_\alpha$，以及点 $z_0:=z_{\alpha\beta}$、$z_\infty:=z_{\beta\alpha}$。
给定 $\varepsilon>0$，按命题~\ref{prop:4-7-2} 选取 $\delta>0$ 和 $\nu_0$，并令
$r:=\delta/2$。若 $\nu\ge\nu_0$、$z\in A_r^\nu(\alpha,\beta)$，且
$\zeta:=(\phi_\alpha^\nu)^{-1}(z)$，则
\[
  d(\zeta,z_0)+d\bigl((\psi^\nu)^{-1}(\zeta),z_\infty\bigr)<2r=\delta.
\]
由于 $z_0=z_{\alpha\beta}$，这里的中心值应为
$u_\alpha(z_{\alpha\beta})$
\sourcefootnote{原书此处把\BookRecordIndex{T0191}节点写成不存在的 $u_{\alpha\beta}$，从而出现
$u_\alpha(u_{\alpha\beta})$；由本段定义 $z_0:=z_{\alpha\beta}$ 以及
命题~\ref{prop:4-7-2} 的结论，唯一有类型的表达式是
$u_\alpha(z_{\alpha\beta})$。}，因而
\[
  d\bigl(u^\nu(z),u_\alpha(z_{\alpha\beta})\bigr)
  =d\bigl(u_\alpha^\nu(\zeta),u_\alpha(z_0)\bigr)<\varepsilon.
\]
式~\eqref{eq:5-2-3} 得证。

现在选取小于 $M$ 的\BookRecordIndex{T0093}单射半径一半的 $\varepsilon>0$，并固定使
式~\eqref{eq:5-2-3} 成立的 $r>0$ 与 $\nu_0\in\mathbb N$。必要时缩小 $r$，可以假设
对每个 $\alpha$，所有满足 $\alpha E\beta$ 的闭圆盘
$\overline B_r(z_{\alpha\beta})$ 两两不交。定义
\[
  S_\alpha:=S^2\setminus
  \bigcup_{\alpha E\beta}\overline B_r(z_{\alpha\beta}).
\]
对每条边 $\alpha E\beta$，用一个标准\BookRecordIndex{T0170}环域
$A_{\alpha\beta}\cong[0,1]\times S^1$ 连接
$\partial B_r(z_{\alpha\beta})\subset S_\alpha$ 与
$\partial B_r(z_{\beta\alpha})\subset S_\beta$，从而构造曲面 $\Sigma_r$。于是，
$A_{\alpha\beta}$ 中每条弧 $[0,1]\times e^{i\theta}$ 都把
$\partial B_r(z_{\alpha\beta})\subset S_\alpha$ 中的一点连接到
$\partial B_r(z_{\beta\alpha})\subset S_\beta$ 中的一点。曲面 $\Sigma_r$ 与 $S^2$
同胚。定义 $u_r:\Sigma_r\to M$，使
\[
  u_r|_{S_\alpha}:=u_\alpha,
\]
并在每个\BookRecordIndex{T0170}环域上，用连接两侧边界对应点之给定像的唯一极短测地线定义 $u_r$。
对 $\nu\ge\nu_0$，类似地定义 $u_r^\nu:\Sigma_r\to M$，使
\[
  u_r^\nu|_{S_\alpha}:=u_\alpha^\nu
  =u^\nu\circ\phi_\alpha^\nu,
\]
并在\BookRecordIndex{T0170}环域上仍用 $M$ 中的测地线连接两侧边界对应点的像。于是，$u_r$ 是各
$u_\alpha$ 的\BookRecordIndex{T0227}连通和，而 $u_r^\nu$ 一致收敛到 $u_r$。

我们断言，当 $\nu$ 充分大时，存在一个同胚
\[
  \psi_r^\nu:\Sigma_r\to S^2,
\]
使得 $u_r^\nu$ 与 $u^\nu\circ\psi_r^\nu$ 同伦。$\psi_r^\nu$ 在 $S_\alpha$ 上的限制
由 $\phi_\alpha^\nu$ 给出，其像为
\[
  S_\alpha^\nu:=\psi_r^\nu(S_\alpha)
  =\phi_\alpha^\nu(S_\alpha)
  =S^2\setminus\bigcup_{\alpha E\beta}
    \phi_\alpha^\nu\bigl(B_r(z_{\alpha\beta})\bigr).
\]
补集 $S^2\setminus\bigcup_\alpha S_\alpha^\nu$ 恰是各\BookRecordIndex{T0170}环域
$A_r^\nu(\alpha,\beta)$ 的并。因此，$\psi_r^\nu$ 可以延拓到\BookRecordIndex{T0170}环域上，成为从
$\Sigma_r$ 到 $S^2$ 的保定向同胚。复合 $u^\nu\circ\psi_r^\nu$ 在每个
$S_\alpha$ 上与 $u_r^\nu$ 相同，并且由式~\eqref{eq:5-2-3}，它在每个\BookRecordIndex{T0170}环域上都与
$u_r^\nu$ 在 $C^0$ 意义下接近。因此 $u^\nu\circ\psi_r^\nu$ 与 $u_r^\nu$ 同伦。
这证明了 (ii)。

下面证明 (iii)。设 $i\in\{1,\ldots,n\}$。我们要证明，对每个顶点 $\alpha$，都有
\[
  z_{\alpha i}=\lim_{\nu\to\infty}(\phi_\alpha^\nu)^{-1}(z_i^\nu).
\]
对从 $\alpha_i$ 到 $\alpha$ 的边数 $k$ 作归纳。当 $k=0$ 时，$\alpha=\alpha_i$，
由式~\eqref{eq:5-1-6}，$z_{\alpha i}=z_{\alpha_i i}:=z_i$，所以结论正是
定义~\ref{def:5-2-1} 的（\BookRecordIndex{T0066}标记点）公理。现在设 $k>0$，并归纳假设：对所有能由不超过
$k-1$ 条边的链从 $\alpha_i$ 到达的顶点，结论都已成立。设 $\alpha$ 可由一条顶点链
$\gamma_0,\ldots,\gamma_k$ 从 $\alpha_i$ 到达，其中
$\gamma_0=\alpha_i$、$\gamma_k=\alpha$（见第 D.2 节）。归纳假设给出
\begin{equation}\label{eq:5-2-4}
  z_{\gamma_{k-1}i}
  =\lim_{\nu\to\infty}
  (\phi_{\gamma_{k-1}}^\nu)^{-1}(z_i^\nu).
\end{equation}
另一方面，由定义~\ref{def:5-2-1} 的（重标度）公理，映射列
\[
  \phi_{\gamma_k\gamma_{k-1}}^\nu
  =(\phi_{\gamma_k}^\nu)^{-1}\circ\phi_{\gamma_{k-1}}^\nu
\]
在 $S^2\setminus\{z_{\gamma_{k-1}\gamma_k}\}$ 上\BookRecordIndex{T0416}逐紧一致收敛到
\[
  z_{\gamma_k\gamma_{k-1}}=z_{\gamma_k i}.
\]
因为 $z_{\gamma_{k-1}i}\ne z_{\gamma_{k-1}\gamma_k}$，这说明
\[
  \begin{aligned}
  z_{\gamma_k i}
  &=\lim_{\nu\to\infty}
    \phi_{\gamma_k\gamma_{k-1}}^\nu(z_{\gamma_{k-1}i})\\
  &=\lim_{\nu\to\infty}
    \phi_{\gamma_k\gamma_{k-1}}^\nu
    \bigl((\phi_{\gamma_{k-1}}^\nu)^{-1}(z_i^\nu)\bigr)\\
  &=\lim_{\nu\to\infty}(\phi_{\gamma_k}^\nu)^{-1}(z_i^\nu).
  \end{aligned}
\]
第二个等号用到了式~\eqref{eq:5-2-4}。因此 (iii) 得证。

断言 (iv) 由 (iii) 和\BookRecordIndex{T0185}交比在 \BookRecordIndex{T0040}Möbius 变换下的不变性立即得到；参见附录 D。
定理~\ref{thm:5-2-2} 得证。
\end{proof}

\begin{remark}\label{rem:5-2-3}
在复几何中，复维 $n$ 的流形上秩为 $k$ 的向量丛的全纯截面，与这些截面的零点集所形成
的复维 $n-k$ 次簇之间存在一种对偶。当 $n-k=1$ 时，这些次簇是曲线；它们还可描述为
全纯映射 $u:\Sigma\to M$ 的像。在许多情形下，全纯曲线的带参数映射描述与不带参数的
次簇描述相当接近；$M=\mathbb{CP}^2$ 的情形见习题~\ref{ex:4-2-7}。在另一些情形中
（例如 $M=\mathbb{CP}^1$），若仅把一条全纯曲线与其像等同，就会丢失大部分信息。
事实上，当 $M$ 本身是黎曼曲面时，代数几何学家也把全纯曲线作为映射研究，并用分歧点
和\BookRecordIndex{T0097}单值化数据编码其几何信息。本章的\BookRecordIndex{T0192}紧性结果所关注的是映射的极限行为，并同时从定义域
和像两方面解释这种行为。\personlater{geluomofu}{格罗莫夫}{Gromov}在其原始论文
\srcite{160} 中以颇为不同的方式理解这种收敛：他从像次簇上的诱导度量来分析它。
\end{remark}
